\begin{proof}[Proof of \cref{obj:prop:regime-and-sanity}]
\begin{enumerate}
\item Fix \(n\ge2\) and \(0<\epsilon\le1\), and define the three benchmark terms by
\[
r_{\mathrm{samp}}=n^{-1/4},
\qquad
r_{\mathrm{sparse}}=(n^2\epsilon)^{-1/7},
\qquad
r_{\mathrm{dense}}=(n\epsilon)^{-1/2}.
\]
All three terms are positive, and
\[
r(n,\epsilon)
=\min\{1,\max\{r_{\mathrm{samp}},
                    \max\{r_{\mathrm{sparse}},r_{\mathrm{dense}}\}\}\}.
\]
Because logarithms preserve order on positive numbers,
\[
\begin{aligned}
r_{\mathrm{sparse}}\le r_{\mathrm{samp}}
&\iff
-\frac17(2\log n+\log\epsilon)\le-\frac14\log n\\
&\iff \log\epsilon\ge-\frac14\log n
\iff \epsilon\ge n^{-1/4},
\\
r_{\mathrm{dense}}\le r_{\mathrm{samp}}
&\iff
-\frac12(\log n+\log\epsilon)\le-\frac14\log n\\
&\iff \log\epsilon\ge-\frac12\log n
\iff \epsilon\ge n^{-1/2}.
\end{aligned}
\]
If \(n^{-1/4}\le\epsilon\), then
\[
n^{-1/2}\le n^{-1/4}\le\epsilon,
\qquad
r_{\mathrm{sparse}}\le r_{\mathrm{samp}},
\qquad
r_{\mathrm{dense}}\le r_{\mathrm{samp}}.
\]
Here the first comparison follows from \(n\ge1\) and
\(-1/2\le-1/4\). Also \(r_{\mathrm{samp}}\le1\), since its exponent is nonpositive. Consequently,
\[
r(n,\epsilon)
=\min\{1,r_{\mathrm{samp}}\}
=r_{\mathrm{samp}}
=n^{-1/4}.
\]
% lean: CausalSmith.Stat.PrivateCateRoughdesign.rate_sampling_regime

\item Suppose that \(n^{-3/5}\le\epsilon\le n^{-1/4}\).
The upper budget bound gives
\[
\log\epsilon\le-\frac14\log n,
\qquad
-\frac14\log n
\le-\frac17(2\log n+\log\epsilon),
\]
and hence
\[
r_{\mathrm{samp}}\le r_{\mathrm{sparse}}.
\]
For the other comparison, logarithms give
\[
\begin{aligned}
r_{\mathrm{dense}}\le r_{\mathrm{sparse}}
&\iff
-\frac12(\log n+\log\epsilon)
\le-\frac17(2\log n+\log\epsilon)\\
&\iff \log\epsilon\ge-\frac35\log n
\iff \epsilon\ge n^{-3/5}.
\end{aligned}
\]
Thus \(r_{\mathrm{dense}}\le r_{\mathrm{sparse}}\).
Moreover, monotonicity in the exponent for a base at least one yields
\[
n^{-1}\le n^{-3/5}\le\epsilon.
\]
Multiplying by \(n>0\), and then using \(n\ge1\), gives
\[
1\le n\epsilon\le n^2\epsilon.
\]
Since the exponent \(-1/7\) is nonpositive,
\[
r_{\mathrm{sparse}}=(n^2\epsilon)^{-1/7}\le1.
\]
The maximum and the cap therefore simplify to
\[
r(n,\epsilon)
=\min\{1,r_{\mathrm{sparse}}\}
=r_{\mathrm{sparse}}
=(n^2\epsilon)^{-1/7}.
\]
% lean: CausalSmith.Stat.PrivateCateRoughdesign.rate_sparse_regime

\item Suppose that \(n^{-1}\le\epsilon\le n^{-3/5}\).
Reversing the preceding logarithmic comparison gives
\[
\begin{aligned}
r_{\mathrm{sparse}}\le r_{\mathrm{dense}}
&\iff
-\frac17(2\log n+\log\epsilon)
\le-\frac12(\log n+\log\epsilon)\\
&\iff \log\epsilon\le-\frac35\log n
\iff \epsilon\le n^{-3/5}.
\end{aligned}
\]
Also,
\[
\epsilon\le n^{-3/5}\le n^{-1/2},
\]
where the second inequality follows from \(n\ge1\).
Taking logarithms and rearranging gives
\[
\log\epsilon\le-\frac12\log n,
\qquad
-\frac14\log n
\le-\frac12(\log n+\log\epsilon),
\]
so that
\[
r_{\mathrm{samp}}\le r_{\mathrm{dense}}.
\]
The lower budget bound gives
\[
n\epsilon\ge1,
\qquad
r_{\mathrm{dense}}=(n\epsilon)^{-1/2}\le1.
\]
It follows that
\[
r(n,\epsilon)
=\min\{1,r_{\mathrm{dense}}\}
=r_{\mathrm{dense}}
=(n\epsilon)^{-1/2}.
\]
% lean: CausalSmith.Stat.PrivateCateRoughdesign.rate_dense_regime

\item If \(\epsilon\le n^{-1}\), then positivity and multiplication by \(n\) give
\[
0<n\epsilon\le1.
\]
A nonpositive power reverses order on positive bases, so
\[
1\le(n\epsilon)^{-1/2}
=r_{\mathrm{dense}}
\le\max\{r_{\mathrm{samp}},
             \max\{r_{\mathrm{sparse}},r_{\mathrm{dense}}\}\}.
\]
Therefore the cap is active:
\[
r(n,\epsilon)=1.
\]
% lean: CausalSmith.Stat.PrivateCateRoughdesign.rate_capped_regime

\item Since \(n\ge1\),
\[
n^{-1/4}\le1.
\]
Applying the sampling-regime identity at \(\epsilon=1\) yields
\[
r(n,1)=n^{-1/4}.
\]
% lean: CausalSmith.Stat.PrivateCateRoughdesign.rate_unit_budget

\item Fix a budget sequence satisfying \(0<\epsilon_n\le1\) for every \(n\ge2\).
The replacement-copy variance inequality
\citep[Theorem 3.1, printed p.~54]{BoucheronLugosiMassart2013},
the bounded-differences exponential-moment bound
\citep[Theorem 6.2, displayed log-mgf conclusion in its proof, printed p.~167]{BoucheronLugosiMassart2013},
and Cayley's labeled-tree formula
\citep[Section 5.6, Theorem 5.40]{KellerTrotterAppliedCombinatorics}
supply the variance, concentration, and counting inputs used in
\cref{obj:thm:matched-risk-frontier}. Applying that result at each \(n\ge2\) gives
\[
2^{-20}r(n,\epsilon_n)
\le R_{n,\epsilon_n}
\le2^{18}r(n,\epsilon_n).
\]
The finite real bounds are regarded here as elements of \([0,\infty]\).
We first establish the numerical equivalence
\[
r(n,\epsilon_n)\longrightarrow0
\quad\Longleftrightarrow\quad
n\epsilon_n\longrightarrow+\infty.
\]
% lean: CausalSmith.Stat.PrivateCateRoughdesign.matched_risk_frontier

\item Suppose that \(r(n,\epsilon_n)\to0\).
For every \(n\ge2\), the dense term is bounded by the uncapped maximum, and monotonicity of the minimum therefore gives
\[
\min\{1,(n\epsilon_n)^{-1/2}\}
\le r(n,\epsilon_n).
\]
For an arbitrary real threshold, set
\[
K_{\mathrm{test}}\in\mathbb R,
\qquad
B_{\mathrm{test}}=\max\{1,K_{\mathrm{test}}\}.
\]
Then
\[
B_{\mathrm{test}}\ge1,
\qquad
\min\{1,B_{\mathrm{test}}^{-1/2}\}>0.
\]
Convergence of the benchmark implies that, eventually,
\[
\min\{1,(n\epsilon_n)^{-1/2}\}
\le r(n,\epsilon_n)
<\min\{1,B_{\mathrm{test}}^{-1/2}\}.
\]
This forces
\[
(n\epsilon_n)^{-1/2}<B_{\mathrm{test}}^{-1/2}.
\]
Indeed, the opposite weak inequality would, by monotonicity of the minimum, imply
\[
\min\{1,B_{\mathrm{test}}^{-1/2}\}
\le\min\{1,(n\epsilon_n)^{-1/2}\},
\]
contradicting the strict bound above.
Both bases are positive, and the power \(-1/2\) is strictly decreasing on positive bases. Hence, eventually,
\[
n\epsilon_n>B_{\mathrm{test}}\ge K_{\mathrm{test}}.
\]
Since the threshold was arbitrary,
\[
n\epsilon_n\longrightarrow+\infty.
\]
% lean: CausalSmith.Stat.PrivateCateRoughdesign.budget_diverges_of_rate_tendsto_zero

\item Conversely, suppose that \(n\epsilon_n\to+\infty\).
First,
\[
n^{-1/4}\longrightarrow0.
\]
For \(n\ge2\), positivity of the budget and \(n\le n^2\) give
\[
n\epsilon_n\le n^2\epsilon_n.
\]
Thus
\[
n^2\epsilon_n\longrightarrow+\infty.
\]
Negative powers tend to zero as their positive bases tend to infinity, so
\[
(n^2\epsilon_n)^{-1/7}\longrightarrow0,
\qquad
(n\epsilon_n)^{-1/2}\longrightarrow0.
\]
Continuity of the finite maximum and minimum now gives
\[
\begin{aligned}
r(n,\epsilon_n)
&=\min\left\{1,\max\left\{
n^{-1/4},
\max\{(n^2\epsilon_n)^{-1/7},(n\epsilon_n)^{-1/2}\}
\right\}\right\}\\
&\longrightarrow
\min\{1,\max\{0,\max\{0,0\}\}\}=0.
\end{aligned}
\]
This completes the numerical equivalence.
% lean: CausalSmith.Stat.PrivateCateRoughdesign.rate_tendsto_zero_of_budget_diverges

\item It remains to transfer convergence between the benchmark and the risk.
Suppose first that \(R_{n,\epsilon_n}\to0\) in \([0,\infty]\).
For \(n\ge2\), positivity of the sampling term and of the cap gives
\[
0<r(n,\epsilon_n).
\]
Fix a tolerance
\[
\eta\in(0,\infty).
\]
Since \(2^{-20}\eta>0\), convergence of the risk implies that, eventually,
\[
R_{n,\epsilon_n}<2^{-20}\eta.
\]
Combining this with the lower frontier bound gives
\[
2^{-20}r(n,\epsilon_n)
\le R_{n,\epsilon_n}
<2^{-20}\eta.
\]
Cancellation of the positive constant yields
\[
0<r(n,\epsilon_n)<\eta.
\]
Therefore \(r(n,\epsilon_n)\to0\), and the numerical equivalence implies
\(n\epsilon_n\to+\infty\).

Conversely, if \(n\epsilon_n\to+\infty\), the numerical equivalence gives
\(r(n,\epsilon_n)\to0\). Consequently,
\[
2^{18}r(n,\epsilon_n)\longrightarrow0
\]
also when these finite values are viewed in \([0,\infty]\).
The upper frontier bound and nonnegativity give, for every \(n\ge2\),
\[
0\le R_{n,\epsilon_n}\le2^{18}r(n,\epsilon_n).
\]
The squeeze principle in \([0,\infty]\) yields
\[
R_{n,\epsilon_n}\longrightarrow0.
\]
Together these implications prove the asserted consistency equivalence.
% lean: CausalSmith.Stat.PrivateCateRoughdesign.private_risk_consistency_iff
\end{enumerate}
\end{proof}
